% Cote 119: offprints, reports, courses, CV and letters (typed and printed part): archive cote 119, site page 282 (PDF page 283).
% typed letter. Draft compared: none.
% AI-transcribed by Claude Opus 5.5 (claude-opus-5-5, adaptive thinking, effort high), Message Batches API, batch r119_full (msgbatch_01DwZNm5Y7YsKBWyAN4sFs8w), 2026-11-13; not yet checked by eye.
\scanpage{119/282}
\inserted{Prof A Grothendieck}
\inserted{\unsure{Les Aumettes|Les Aumnettes}}
\inserted{8470 Mormoiron}
\inserted{FRANCE}

Dear Alexander,
16.04.91
This is a letter which goes backwards, because I started it a long time
ago, and it never got sent.  Then today I got your letter (or life sign, as
you say), which delights us both with your gift of language and direct
reaction!  It was Margaret who suggested you would enjoy a card from Iona,
so she was especially pleased.
I would love to hear more of your `vacations' in mathematics; and there
are many all over the world who would also love to hear of your new vision,
or even to have some inkling of it, without necessarily following you `to
the beyond'!
We were also intrigued by your hint of a `near end to this pursuit'.  I
do hope you are not ill in any way (the suggestion is belied by what else
you write).  Are you expecting in the near future some other kind of general
end?
Of the mathematical things you have told me over some years, perhaps
the one that most baffles and intrigues me, because it is so directly
related to some of my ideas, is the area of the Teichm\"uller groupoid, and
its presentations.  It agrees with the thought that groups are inadequate to
express our intuitions of symmetry, and we must pass to groupoids, and other
more complex structures.  But I still cannot recover your results.
The theme has lead on to ideas of `higher order symmetry', for which
there is a lot to do.  The line I have been taking is that multiple
groupoids are a clear generalisation of groups, and so their study has a
clear relation to `higher order symmetry', whatever that may be.  One aspect
is that an object reflecting higher order symmetry should be the
coefficients for a cohomology describing extensions, and crossed modules.
The simplest example is that of groups, where the `automorphism structure'
of a group $G$ is really the crossed module $G \longrightarrow \mathrm{Aut}\, G$ , and it is this
crossed module which is the coefficients for extensions, as realised by
Dedecker.  At the next stage, the automorphisms of a crossed module form
part of a structure which may be described as a crossed square, or a
simplicial group whose Moore complex is of length 2, or what Gilbert and I
called a `braided crossed module' (over a groupoid).  But conversations with
Larry Breen suggest that even at level three the symmetry of a bitorseur
will require equivalences rather than isomorphisms, and so the dreaded beast
of structures up to coherence emerges again.
06.05.91
I should also mention the work of my student John Shrimpton.  In order
to interpret if possible what you said about the Teichm\"uller groupoid, we
decided to consider the similar problems for graphs.  Is it possible to give
a presentation of an `automorphism groupoid' (whatever that is) of a graph
from the way it is built up out of simpler graphs?  This was expected to be
analogous to the Van Kampen Theorem.
We were unable to do this, but we had an interesting answer to
describing the symmetry of a graph by a `group-graph'.  A directed graph
with base loops is just a 1-truncated simplicial set, and so is a functor
category and hence a topos.  So the automorphism object of such a graph is a
1-truncated simplicial group.  As such, it has a notion of $\pi_1$ and $\pi_2$ : the
fundamental group is the group of \textit{outer automorphisms, and} $\pi_2$ is the
\textit{centre} of the original graph (it is a module over $\pi_1$ ).  John defines an
\textit{inner subgraph} I of a graph $\Gamma$ to be maximal with respect to the properties
of being complete (all edge sets I(x,y) are bijective), full, and
symmetrically embedded (any automorphism of I extends to an automorphism of
$\Gamma$ which is the identity on $\Gamma\backslash$I).  It turns out that inner automorphisms of
$\Gamma$ are precisely those which restrict to an automorphism of each inner
subgraph.  John also shows that in the finite case, $\pi_2$ has every element

\mspage{1}

\margin{282}
